\documentclass[12pt,a4paper]{ctexart}

\usepackage[margin=2.2cm]{geometry}
\usepackage{amsmath,amssymb,amsfonts}
\usepackage{algorithm}
\usepackage{algpseudocode}
\usepackage{graphicx}
\usepackage{textcomp}
\usepackage{xcolor}
\usepackage{multirow}
\usepackage{booktabs}
\usepackage{float}
\usepackage{hyperref}
\usepackage{cite}

\algrenewcommand\algorithmicrequire{\textbf{输入：}}
\algrenewcommand\algorithmicensure{\textbf{输出：}}

\title{基于双层对抗训练的网络威胁检测方法研究}
\author{朱军}
\date{\today}

\begin{document}

\maketitle

\begin{abstract}
网络入侵检测系统（IDS）面临的核心挑战在于：针对固定数据集训练的静态检测器在面对自适应攻击者时性能急剧下降，而现有训练方法大多假设威胁环境是被动且静止的。本文提出\textbf{双层对抗强化学习（Bi-ARL）}框架，将攻防交互形式化为 Stackelberg 博弈，并在每次外层防御者更新前通过内层循环将攻击者迭代至近似最优响应，从而实现比同步更新 MARL 更强的鲁棒性训练目标。在此基础上，本文将相同的双层对抗训练原则迁移至更强的监督式表格检测器，构建\textbf{BiAT 框架}（BiAT-MLP 与 BiAT-FTTransformer）。在 NSL-KDD、UNSW-NB15 与 CIC-IDS2017 三个公开基准上的系统评估表明：（1）Bi-ARL 在 NSL-KDD 上以 80.17\% 的平均 F1 居强化学习基线之首，相较于 Vanilla PPO 将误报率降低 29.16 个百分点（统计显著，$p<0.05$）；（2）BiAT-FTTransformer 在 UNSW-NB15 上取得 89.51\% 的平均 F1 和 14.97\% 的平均 FPR，在所有神经网络检测器中取得最低误报率；（3）FGSM 与 PGD 鲁棒性分析表明双层训练在小扰动预算下提供部分鲁棒性，同时揭示了内层循环收敛质量是当前主要瓶颈。本文的研究建立了显式 leader-follower 双层优化作为对抗鲁棒 IDS 训练的有效方法论，并明确指出了进一步提升的工程路径。
\end{abstract}

\noindent\textbf{关键词：}网络威胁检测；双层优化；Stackelberg 博弈；对抗训练；强化学习；入侵检测系统

%=============================================================================
\section{引言}
%=============================================================================

网络安全仍是现代计算基础设施的核心挑战。传统 IDS 依赖签名或静态统计规则，在面对零日攻击、分布漂移与对抗规避时固有地脆弱。近年来，基于学习的 IDS 研究已推进到自监督泛化、开放集识别以及基于 Transformer 的特征建模等方向\cite{liao2024survey,golchin2024ssclids,xi2024idsmtran}。然而，一个更深层的问题仍未得到充分解决：大多数 IDS 训练流程假设威胁环境是\emph{被动的}——攻击者不会针对已部署的模型进行自适应优化。经验证据表明，在单一数据集上表现优异的分类器，在跨数据集测试或针对性扰动下可能急剧退化\cite{cantone2024crossdataset,wang2023robustnessnids}，这正是因为它们从未针对具有策略性的自适应对手进行训练。

强化学习为自适应防御提供了自然框架，然而现有 RL-based IDS 管线大多仍采用静态扰动模型或\emph{同步}更新的对手进行训练。同步更新产生的是共适应，而非真正的最坏情形强化：双方均未针对对方的最优响应进行优化。这一区别至关重要：Stackelberg 博弈形式化——防御者提前承诺策略、攻击者据此最优响应——比普通多智能体联合训练提供了更强的鲁棒性目标。

本文将 Stackelberg 形式化付诸实践。我们提出\textbf{双层对抗强化学习（Bi-ARL）}，其中外层防御者更新以内层攻击者近似最优响应为条件。随后，考虑到 PPO 骨干网络作为表格检测器相对较弱，我们将同样的双层对抗训练原则迁移至现代监督式检测器，形成\textbf{BiAT 框架}：BiAT-MLP 与 BiAT-FTTransformer。

\subsection{主要贡献}

\begin{enumerate}
    \item \textbf{Bi-ARL 框架。} 将入侵检测训练形式化为双层 Stackelberg 优化问题，实现了基于 PPO 的实用管线，并采用基于 KL 散度的内层收敛准则。

    \item \textbf{NSL-KDD 实证验证。} Bi-ARL 达到 80.17\% 的平均 F1，并将 FPR 相较 Vanilla PPO 降低 29.16 个百分点（统计显著），证实显式 leader-follower 训练比同步更新 MARL 产生更优的精确率-召回率权衡。

    \item \textbf{BiAT 框架：监督式双层对抗训练。} 将核心双层思想迁移至 BiAT-MLP 和 BiAT-FTTransformer，在保留 Stackelberg 训练原则的同时引入更强的表格建模骨干。

    \item \textbf{UNSW-NB15 最优神经网络性能。} BiAT-FTTransformer 在 UNSW-NB15 上取得 89.51\% 的 F1 与 14.97\% 的 FPR，在所有神经网络检测器中取得最低误报率，与 LightGBM 的 F1 差距缩小至 2.35 个百分点。

    \item \textbf{全面的对抗鲁棒性研究。} 对所有模型族（含BiAT 框架）进行 FGSM 与 PGD 鲁棒性评估，辅以统计显著性检验与超参数敏感性分析，对框架的能力与局限性提供诚实的定量表征。
\end{enumerate}

%=============================================================================
\section{相关工作}
%=============================================================================

\subsection{入侵检测系统研究}

传统 IDS 依赖签名匹配，对已知攻击模板有效，但对零日攻击或快速演化威胁效果有限。学习式 IDS 已覆盖传统机器学习、深度神经网络、图模型及 IoT 场景检测\cite{liao2024survey,moustafa2023explainable}。与此同时，数据集质量已成为关键瓶颈：许多公开 NIDS 基准包含过时流量模式、快捷特征或过于宽松的划分假设，可能显著高估部署性能\cite{ring2019nids,datasetreview2025}。这与本文直接相关——Bi-ARL 在 NSL-KDD、UNSW-NB15 和 CIC-IDS2017 上的结论存在明显差异，与 Cantone 等人记录的跨数据集不稳定性一致\cite{cantone2024crossdataset}。

\subsection{泛化能力与现代深度 IDS}

超越单一数据集封闭集评测是近年的重要趋势。SSCL-IDS 通过自监督对比学习提升未知流量分布下的检测泛化能力\cite{golchin2024ssclids}；GNN-SSL 将自监督学习扩展至基于图神经网络的流量建模\cite{xu2024gnnssl}；IDS-MTran 采用多尺度 Transformer 结构并在 NSL-KDD 和 UNSW-NB15 上表现强劲\cite{xi2024idsmtran}；2025 年的跨数据集 Transformer-IDS 研究进一步引入校准与 AUC 优化\cite{xin2025crossdatasettransformer}。本文 BiAT-FTTransformer 所使用的 FT-Transformer 主干由 Gorishniy 等人提出，专为表格数据设计\cite{gorishniy2021fttransformer}。这些工作构成了远超传统 RL-only 基线的现代对比集，但它们均未将自适应攻击者显式纳入训练过程。

\subsection{鲁棒性、开放集检测与跨数据集评估}

Wang 等人指出对抗扰动与分布漂移应作为机器学习 NIDS 的一级评测标准\cite{wang2023robustnessnids}。Apruzzese 等人系统建模了针对网络分类器的现实对抗攻击，表明精心设计的扰动会导致显著退化\cite{apruzzese2022adversarial_nids}。VAEMax 通过结合 OpenMax 与变分自编码器探索了开放集入侵检测\cite{qiu2024vaemax}。SHIELD/NF-UQ-NIDS 与 5G-NIDD 等新数据资源的出现，预示着领域正在向更标准化、当代化的流量数据集迈进\cite{shield2025,sarhan2021nfuqnids,siriwardhana2025nidd}。本文的鲁棒性分析直接继承这一方向，对 RL-based 与监督式双层模型均进行了 FGSM 和 PGD 评估。

\subsection{对抗机器学习基础}

Goodfellow 等人形式化了对抗样本威胁\cite{goodfellow2015fgsm}，Madry 等人将其扩展为迭代 PGD 攻击\cite{madry2018pgd}，为本文BiAT 框架模型的训练内循环提供了理论基础。Biggio 等人较早系统阐释了测试时逃逸攻击\cite{biggio2013evasion}，Yuan 等人的综述进一步总结了深度学习对抗攻击与防御的发展脉络\cite{yuan2019adv_survey}。Carlini 和 Wagner 证明许多针对 FGSM 的防御在更强攻击下仍失效\cite{carlini2017cw}。Zhang 等人通过 TRADES 刻画了鲁棒性-准确性权衡\cite{zhang2019trades}——本文BiAT 框架 的混合损失 $L = (1-\lambda)L_{\text{clean}} + \lambda L_{\text{adv}}$ 正是这一权衡的直接实例。

\subsection{对抗强化学习与 Stackelberg 博弈}

将 RL 应用于安全博弈可自然建模顺序攻防交互\cite{anderson2016rl_security}。Nash-Q、分布式入侵响应 RL、MARL 与 MADDPG 等标准多智能体 RL 方案通常同步更新双方，产生共适应而非真正的最坏情形强化\cite{hu2003nashq,malialis2015multiagent,busoniu2008marl,lowe2017maddpg}。Stackelberg 形式化通过指定一方为先行者（领导者）、另一方据此最优响应（跟随者）来解决这一问题\cite{vonStackelberg1952,colson2007bilevel,sinha2018bilevel,bard1998bilevel}。双层优化为实现这一结构提供了数学工具。Pinto 等人在物理控制中应用了类似结构\cite{pinto2017robust_rl}；Li 等人将对抗 RL 扩展至软件定义网络\cite{li2022advrl_security}。本文 Bi-ARL 显式实现了收敛检测的双层 Stackelberg 结构，区别于同步更新 MARL 基线。

%=============================================================================
\section{方法}
%=============================================================================

本节介绍 Bi-ARL 框架与BiAT 框架 扩展。首先给出正式威胁模型，然后描述 Stackelberg 博弈形式化、双层优化算法及监督式双层迁移。

\subsection{威胁模型}
\label{sec:threatmodel_cn}

我们在训练时假设\textbf{白盒自适应对手}，测试时在白盒和梯度攻击两种条件下进行评估。

\textbf{攻击者能力。} 攻击者完全了解防御者的模型结构与参数，可对输入特征计算梯度，目标是构造满足 $\|\delta\|_\infty \leq \epsilon$ 的扰动以最大化检测器的漏报率（FNR）。类别型特征（如协议类型）被排除在扰动范围外。

\textbf{防御者目标。} 防御者寻求一个检测策略 $\pi_{\mathcal{D}}$，即使在针对攻击者最优响应的条件下，也能最小化综合漏报与误报的损失。

\textbf{训练与部署的区分。} 在\emph{训练}阶段，攻击者为白盒并被反复更新；在\emph{部署}（测试）阶段，采用多种扰动预算的 FGSM 与 PGD 攻击衡量已训练模型的实际鲁棒性。

\subsection{问题形式化}

设 $\mathcal{S}$ 为表格化流量特征的状态空间，维数取决于数据集（NSL-KDD 为 41 维，UNSW-NB15 为 42 维，CIC-IDS2017 预处理后为 78 维）。环境包含两个智能体：攻击者（$\mathcal{A}$）与防御者（$\mathcal{D}$）。

本文将该问题建模为 \textbf{Stackelberg 博弈}，其外层-内层结构同时可视为一种面向安全训练的鲁棒优化近似\cite{ben2009robust}：

\begin{equation}
\max_{\theta_{\mathcal{D}}} J(\theta_{\mathcal{D}}, \theta_{\mathcal{A}}^*) =
\mathbb{E}_{\tau \sim \pi_{\mathcal{D}}, \pi_{\mathcal{A}}^*}
\left[\sum_{t=0}^{T} \gamma^t R_{\mathcal{D}}(s_t, a_t^{\mathcal{D}}, a_t^{\mathcal{A}*})\right]
\label{eq:outer_cn}
\end{equation}

\begin{equation}
\text{s.t.} \quad
\theta_{\mathcal{A}}^* =
\arg\max_{\theta_{\mathcal{A}}}
\mathbb{E}_{\tau \sim \pi_{\mathcal{A}}}
\left[\sum_{t=0}^{T} \gamma^t R_{\mathcal{A}}(s_t, a_t^{\mathcal{D}}, a_t^{\mathcal{A}})\right]
\label{eq:inner_cn}
\end{equation}

其中 $\gamma = 0.99$ 为折扣因子，$\tau$ 表示一条轨迹。式~\eqref{eq:outer_cn} 为外层（防御者）问题，式~\eqref{eq:inner_cn} 为内层（攻击者最优响应）问题。这种嵌套结构正是 Bi-ARL 与同步更新 MARL 的本质区别。

\subsection{奖励函数}

\subsubsection{攻击者奖励}

\begin{equation}
R_{\mathcal{A}}=
\begin{cases}
+2.0, & \text{漏报（攻击规避检测）}\\
-1.0, & \text{检出（攻击被检测到）}
\end{cases}
\end{equation}

\subsubsection{防御者奖励}

\begin{equation}
R_{\mathcal{D}}=
\begin{cases}
+1.0, & \text{真正例（TP）}\\
+0.5, & \text{真负例（TN）}\\
-1.0, & \text{假正例（FP）}\\
-2.0, & \text{假负例（FN）}
\end{cases}
\end{equation}

非对称惩罚（$|r_{\text{FN}}| > |r_{\text{TP}}|$）反映了漏报攻击在运营层面的更高代价。对于 UNSW-NB15 和 CIC-IDS2017，假正例惩罚增大至 $-2.5$ 以在高类别不平衡下稳定训练。

\subsection{基于 PPO 的策略优化}

双方智能体均使用 PPO（近端策略优化）进行更新\cite{schulman2017ppo}，并结合 GAE（Generalized Advantage Estimation）稳定优势估计\cite{schulman2015gae}，最大化截断代理目标：

\begin{equation}
L^{\text{CLIP}}(\theta)=
\mathbb{E}_t
\left[
\min\left(
r_t(\theta)\hat{A}_t,\;
\text{clip}(r_t(\theta),1-\epsilon_{\text{clip}},1+\epsilon_{\text{clip}})\hat{A}_t
\right)
\right]
\end{equation}

其中 $r_t(\theta) = \pi_\theta(a_t|s_t)/\pi_{\theta_{\text{old}}}(a_t|s_t)$，$\hat{A}_t$ 为 GAE 优势估计，$\epsilon_{\text{clip}} = 0.2$，$\lambda_{\text{GAE}} = 0.95$。

\subsection{双层优化算法}

\subsubsection{内层：攻击者逼近最优响应}

攻击者在固定防御者策略 $\pi_{\mathcal{D}}^{\text{fixed}}$ 下更新至多 $K_{\text{inner}}$ 步，收敛条件为 KL 散度：

\begin{equation}
D_{\text{KL}}\left(\pi_{\mathcal{A}}^{\text{old}}\|\pi_{\mathcal{A}}^{\text{new}}\right) < \tau_{\text{KL}} = 0.01
\end{equation}

NSL-KDD 采用 $K_{\text{inner}}=5$；在更困难的现代数据集上采用 $K_{\text{inner}}=1$（出于稳定性考虑，见第~\ref{sec:experiments_cn} 节）。

\subsubsection{外层：防御者鲁棒化}

攻击者收敛至近似最优响应 $\theta_{\mathcal{A}}^*$ 后，防御者针对该冻结的最坏情形对手执行一步 PPO 更新。

\begin{algorithm}[H]
\caption{双层对抗强化学习（Bi-ARL）}
\begin{algorithmic}[1]
\Require 训练轮数 $N$，内层最大步数 $K_{\text{inner}}$，KL 阈值 $\tau_{\text{KL}}$
\Ensure 防御者策略 $\pi_{\mathcal{D}}(\cdot|\theta_{\mathcal{D}})$
\State 随机初始化 $\theta_{\mathcal{D}}, \theta_{\mathcal{A}}$
\For{episode $e=1$ 到 $N$}
    \State \textbf{内层：攻击者最优响应}
    \State $\theta_{\mathcal{A}}^{\text{old}}\leftarrow\theta_{\mathcal{A}}$
    \For{$k=1$ 到 $K_{\text{inner}}$}
        \State 采样轨迹 $\tau \sim \pi_{\mathcal{A}}(\cdot|\theta_{\mathcal{A}}),\; \pi_{\mathcal{D}}^{\text{fixed}}$
        \State 更新 $\theta_{\mathcal{A}} \leftarrow \theta_{\mathcal{A}} + \alpha \nabla_{\theta_{\mathcal{A}}} L^{\text{CLIP}}(\theta_{\mathcal{A}})$
        \State 计算 $D_{\text{KL}}(\pi_{\mathcal{A}}^{\text{old}}\|\pi_{\mathcal{A}})$
        \If{$D_{\text{KL}}<\tau_{\text{KL}}$}
            \State \textbf{break} \Comment{攻击者已收敛}
        \EndIf
    \EndFor
    \State $\theta_{\mathcal{A}}^* \leftarrow \theta_{\mathcal{A}}$ \Comment{冻结最优攻击者}
    \State \textbf{外层：防御者鲁棒化}
    \State 采样轨迹 $\tau \sim \pi_{\mathcal{A}}^*(\cdot|\theta_{\mathcal{A}}^*),\; \pi_{\mathcal{D}}(\cdot|\theta_{\mathcal{D}})$
    \State 更新 $\theta_{\mathcal{D}} \leftarrow \theta_{\mathcal{D}} + \beta \nabla_{\theta_{\mathcal{D}}} L^{\text{CLIP}}(\theta_{\mathcal{D}})$
\EndFor
\end{algorithmic}
\end{algorithm}

\subsection{BiAT 框架：双层对抗训练向监督式检测器的迁移}
\label{sec:biat_cn}

Bi-ARL 确立了 Stackelberg 训练原则，但 PPO 策略网络作为表格分类器相对较弱。BiAT 框架 将相同的双层思想迁移至具有更强归纳偏置的监督式检测器。

\textbf{BiAT-MLP}：三层 MLP（$256\to128\to2$，含批归一化与 Dropout），采用 PGD 内层扰动（$\epsilon=0.04$）。

\textbf{BiAT-FTTransformer}：FT-Transformer 将每个数值特征通过学习的线性映射转化为 $d$ 维嵌入，前置 \texttt{[CLS]} 标记后经标准 Transformer 编码层处理。使用 \texttt{[CLS]} 表示进行分类。该结构的逐特征注意力机制特别适合特征交互非平凡的表格 IDS 数据。采用较温和的对抗内层参数（$\epsilon=0.02$，$\alpha=0.005$，2步 PGD）。

两种BiAT 框架模型均最小化：
\begin{equation}
L = (1-\lambda)L_{\text{clean}} + \lambda L_{\text{adv}}
\label{eq:biat_cn}
\end{equation}

其中 $L_{\text{clean}}$ 为干净样本的交叉熵，$L_{\text{adv}}$ 为 PGD 扰动样本的交叉熵，$\lambda$ 控制对抗损失权重。该形式与 TRADES 的鲁棒性-准确性权衡框架一致。

BiAT 框架 提供了清晰的消融链：
\begin{center}
普通 FT-Transformer（仅干净训练）$\to$ BiAT-MLP $\to$ BiAT-FTTransformer
\end{center}

这条消融链将（a）更强骨干网络和（b）双层对抗训练的贡献分离开来。

\subsection{理论分析}

Bi-ARL 近似于 Stackelberg 均衡：防御者针对近似最优响应的攻击者进行优化，而非针对同步共适应的对手。从理论上看，这一训练目标是一个极小极大问题，原则上比标准 ERM 或同步 MARL 训练提供更强的鲁棒性保证。实践中，内层循环以有限步数和 KL 收敛准则运行，因此结果是工程近似而非认证的 Stackelberg 均衡。第~\ref{sec:experiments_cn} 节对该近似与理论理想之间的经验差距进行了定量刻画。

%=============================================================================
\section{实验}
\label{sec:experiments_cn}
%=============================================================================

\subsection{实验设置}

\subsubsection{数据集}

本文在三个公开表格化 IDS 基准上进行评估：
\begin{itemize}
    \item \textbf{NSL-KDD}：\texttt{KDDTrain+}（125,973 条）与 \texttt{KDDTest+}（22,544 条），41 维特征\cite{tavallaee2009nslkdd}。主要用于受控的 RL 消融实验。
    \item \textbf{UNSW-NB15}：官方训练集（175,341 条）与测试集（82,332 条），去除标识与辅助列后 42 维特征\cite{moustafa2015unswnb15}。代表更现代、更有挑战性的基准。
    \item \textbf{CIC-IDS2017}：基于流量的基准，78 维特征\cite{sharafaldin2018cicids2017}。主要设定采用随机分层划分，按工作日划分作为更严格的泛化压力测试。
\end{itemize}

\subsubsection{对比基线}

\begin{enumerate}
    \item Vanilla PPO：无对抗训练的单智能体 RL；
    \item MARL：不采用双层嵌套的同步更新；
    \item 普通 FT-Transformer：仅干净监督训练，作为BiAT 框架 的消融骨干基线；
    \item LSTM-IDS：两层 LSTM 监督式基线\cite{staudemeyer2019lstm}；
    \item HGBT-IDS：直方图梯度提升；
    \item XGBoost-IDS 与 LightGBM-IDS：现代 boosted-tree 基线；
    \item BiAT-MLP 与 BiAT-FTTransformer：BiAT 框架 双层对抗训练监督式检测器。
\end{enumerate}

\subsubsection{训练细节}

除特别说明外，所有结果均报告 4 个随机种子（42、123、3407、8888）的均值与标准差。NSL-KDD 上 RL 训练使用 100 个 episode。UNSW-NB15 上采用 $K_{\text{inner}}=1$ 以稳定训练。LSTM-IDS 训练 20 个 epoch，HGBT-IDS 使用深度 8、迭代 300 次的配置。BiAT-FTTransformer 采用 $d_{\text{token}}=64$、3 层 Transformer 编码器、AdamW（$\text{lr}=7\times10^{-4}$，$\text{wd}=10^{-4}$）。

\subsection{评测指标}

\begin{itemize}
    \item 召回率（Recall，即 TPR）：$TP/(TP+FN)$
    \item 精确率（Precision）：$TP/(TP+FP)$
    \item F1 值：精确率与召回率的调和平均
    \item 误报率（FPR）：$FP/(FP+TN)$——已部署 IDS 的关键运营指标
\end{itemize}

关键对比的统计显著性采用 Wilcoxon 符号秩检验（非参数，适用于 $n=4$ 配对种子）。

\subsection{主结果}

\subsubsection{NSL-KDD}

表~\ref{tab:nsl_cn} 给出 NSL-KDD 主结果。Bi-ARL 在强化学习基线中取得最高平均 F1（80.17\%），相较 Vanilla PPO 将 FPR 降低 29.16 个百分点（Cohen's $d=1.66$，大效应量）。

\begin{table}[H]
\centering
\caption{NSL-KDD 主结果（4 个随机种子均值 $\pm$ 标准差）}
\label{tab:nsl_cn}
\begin{tabular}{lcccc}
\toprule
\textbf{方法} & \textbf{Recall} & \textbf{Precision} & \textbf{F1} & \textbf{FPR} \\
\midrule
Vanilla PPO & $76.39{\pm}19.89$\% & $79.75{\pm}20.65$\% & $74.12{\pm}2.99$\% & $39.57{\pm}46.94$\% \\
MARL & $70.97{\pm}10.91$\% & $91.65{\pm}5.60$\% & $79.32{\pm}4.79$\% & $9.54{\pm}7.73$\% \\
LSTM-IDS & $65.51{\pm}1.64$\% & $\mathbf{96.62{\pm}0.02}$\% & $78.07{\pm}1.18$\% & $3.03{\pm}0.08$\% \\
HGBT-IDS & $67.65{\pm}0.71$\% & $96.80{\pm}0.04$\% & $79.64{\pm}0.50$\% & $2.96{\pm}0.02$\% \\
XGBoost-IDS & $67.10{\pm}0.32$\% & $96.80{\pm}0.03$\% & $79.26{\pm}0.22$\% & $2.93{\pm}0.03$\% \\
LightGBM-IDS & $64.00{\pm}0.24$\% & $96.67{\pm}0.06$\% & $77.01{\pm}0.18$\% & $\mathbf{2.92{\pm}0.04}$\% \\
\midrule
\textbf{Bi-ARL（本文）} & $72.38{\pm}7.46$\% & $90.56{\pm}3.68$\% & $\mathbf{80.17{\pm}3.03}$\% & $10.41{\pm}5.47$\% \\
\bottomrule
\end{tabular}
\end{table}

Bi-ARL 在强化学习基线中取得最优的精确率-召回率平衡，相较 Vanilla PPO 的 FPR 降低具有大效应量（Cohen's $d=1.66$）。Bi-ARL 并非绝对 FPR 最低，因为树模型和 LSTM 在干净数据上 FPR 更低，但代价是未针对自适应攻击者进行训练。

\subsubsection{UNSW-NB15}

表~\ref{tab:unsw_cn} 展示 UNSW-NB15 结果。排名相较 NSL-KDD 发生了显著变化。原始 Bi-ARL 的 RL 骨干表现出高方差和较差的绝对性能，证实基于 PPO 的表格检测器在现代基准上不足。BiAT 框架模型明显更强：BiAT-FTTransformer 相较 BiAT-MLP 将 F1 提升 3.11 个点、FPR 降低 13.34 个点，并在所有评测方法中取得最低平均 FPR。

新增的 Transformer-IDS 基线（仅干净监督训练，无对抗训练）在 UNSW-NB15 上取得 90.57\% 的干净 F1，略高于 BiAT-FTTransformer（89.51\%）。然而，在对抗鲁棒性分析（见第~\ref{sec:rob_cn} 节）中，BiAT-FTTransformer 在 FGSM-$\epsilon$0.3 下 FPR 比 Transformer-IDS 低 16.3 个百分点（54.56\% 对 70.86\%），验证了双层对抗训练以约 1\% 的干净 F1 换取显著鲁棒性提升的权衡关系。

\begin{table}[H]
\centering
\caption{UNSW-NB15 主结果（4 个随机种子均值 $\pm$ 标准差）}
\label{tab:unsw_cn}
\begin{tabular}{lcccc}
\toprule
\textbf{方法} & \textbf{Recall} & \textbf{Precision} & \textbf{F1} & \textbf{FPR} \\
\midrule
Bi-ARL & $66.56{\pm}47.14$\% & $48.80{\pm}35.47$\% & $54.13{\pm}36.12$\% & $53.56{\pm}53.94$\% \\
Transformer-IDS（无对抗训练） & $92.97{\pm}1.43$\% & $88.36{\pm}3.05$\% & $90.57{\pm}0.95$\% & $15.17{\pm}4.74$\% \\
BiAT-MLP & $93.63{\pm}0.54$\% & $80.21{\pm}0.62$\% & $86.40{\pm}0.27$\% & $28.31{\pm}1.22$\% \\
\textbf{BiAT-FTTransformer} & $90.90{\pm}0.81$\% & $88.19{\pm}1.82$\% & $\mathbf{89.51{\pm}0.54}$\% & $\mathbf{14.97{\pm}2.69}$\% \\
LSTM-IDS & $\mathbf{97.44{\pm}1.05}$\% & $79.53{\pm}1.19$\% & $87.57{\pm}0.31$\% & $30.77{\pm}2.57$\% \\
HGBT-IDS & $98.30{\pm}0.27$\% & $82.07{\pm}0.07$\% & $89.45{\pm}0.15$\% & $26.31{\pm}0.08$\% \\
XGBoost-IDS & $97.02{\pm}0.08$\% & $86.17{\pm}0.08$\% & $91.27{\pm}0.03$\% & $19.08{\pm}0.13$\% \\
LightGBM-IDS & $96.83{\pm}0.15$\% & $\mathbf{87.38{\pm}0.07}$\% & $91.86{\pm}0.03$\% & $17.14{\pm}0.13$\% \\
\bottomrule
\end{tabular}
\end{table}

\subsubsection{CIC-IDS2017}

随机分层协议下（表~\ref{tab:cic_cn}），树模型几乎接近饱和（约 99.74\% F1），证实现代表格方法对该数据集的高度可学习性。BiAT 框架 相较旧 RL 主线有明显改进，但仍落后于树模型和 LSTM-IDS。BiAT-FTTransformer 相较 BiAT-MLP 的提升仅 0.36 个 F1 点，表明 CIC-IDS2017 特征在无需强化对抗训练的情况下已具备高可学习性。

\begin{table}[H]
\centering
\caption{CIC-IDS2017 随机分层主结果（4 个随机种子均值 $\pm$ 标准差）}
\label{tab:cic_cn}
\begin{tabular}{lcccc}
\toprule
\textbf{方法} & \textbf{Recall} & \textbf{Precision} & \textbf{F1} & \textbf{FPR} \\
\midrule
BiAT-MLP & $98.65{\pm}0.33$\% & $77.39{\pm}0.41$\% & $86.73{\pm}0.15$\% & $7.16{\pm}0.19$\% \\
BiAT-FTTransformer & $98.20{\pm}0.61$\% & $78.24{\pm}0.55$\% & $87.09{\pm}0.16$\% & $6.78{\pm}0.26$\% \\
LSTM-IDS & $98.15{\pm}0.62$\% & $92.01{\pm}0.74$\% & $94.98{\pm}0.13$\% & $2.12{\pm}0.23$\% \\
HGBT-IDS & $99.85{\pm}0.01$\% & $99.63{\pm}0.01$\% & $99.74{\pm}0.01$\% & $\mathbf{0.09{\pm}0.00}$\% \\
XGBoost-IDS & $\mathbf{99.89{\pm}0.01}$\% & $99.59{\pm}0.01$\% & $99.74{\pm}0.00$\% & $0.10{\pm}0.00$\% \\
LightGBM-IDS & $99.87{\pm}0.01$\% & $\mathbf{99.62{\pm}0.02}$\% & $\mathbf{99.75{\pm}0.01}$\% & $\mathbf{0.09{\pm}0.00}$\% \\
\bottomrule
\end{tabular}
\end{table}

\subsection{消融实验}

NSL-KDD 消融结果（表~\ref{tab:abl_nsl_cn}）表明，去除攻击者（Vanilla PPO）使 FPR 上升 29.16 个百分点，是对抗感知训练驱动 FPR 降低的最直接证据。完整 Bi-ARL 在强化学习变体中取得最高 F1，证实双层嵌套相较同步 MARL 具有增量价值。

在 UNSW-NB15 上（表~\ref{tab:abl_unsw_cn}），消融趋势不如 NSL-KDD 清晰：固定攻击者变体取得比完整 Bi-ARL 更低的 FPR（37.45\% 对 53.56\%），表明内层循环在更困难数据集上尚不稳定，这是未来工作的主要改进方向。

\begin{table}[H]
\centering
\caption{NSL-KDD 消融结果（4 个随机种子均值 $\pm$ 标准差）}
\label{tab:abl_nsl_cn}
\begin{tabular}{lccc}
\toprule
\textbf{变体} & \textbf{Recall} & \textbf{F1} & \textbf{FPR} \\
\midrule
去除内层（固定攻击者）& $70.98{\pm}3.12$\% & $78.78{\pm}2.84$\% & $10.82{\pm}2.31$\% \\
去除攻击者（Vanilla PPO）& $\mathbf{76.39{\pm}19.89}$\% & $74.12{\pm}2.99$\% & $39.57{\pm}46.94$\% \\
去除双层（MARL）& $70.97{\pm}10.91$\% & $79.32{\pm}4.79$\% & $\mathbf{9.54{\pm}7.73}$\% \\
\midrule
\textbf{完整 Bi-ARL} & $72.38{\pm}7.46$\% & $\mathbf{80.17{\pm}3.03}$\% & $10.41{\pm}5.47$\% \\
\bottomrule
\end{tabular}
\end{table}

\begin{table}[H]
\centering
\caption{UNSW-NB15 消融结果（4 个随机种子均值 $\pm$ 标准差）}
\label{tab:abl_unsw_cn}
\begin{tabular}{lccc}
\toprule
\textbf{变体} & \textbf{Recall} & \textbf{F1} & \textbf{FPR} \\
\midrule
去除内层（固定攻击者）& $62.95{\pm}7.43$\% & $56.07{\pm}8.12$\% & $\mathbf{37.45{\pm}9.87}$\% \\
去除攻击者（Vanilla PPO）& $66.21{\pm}29.85$\% & $\mathbf{59.42{\pm}15.08}$\% & $58.72{\pm}31.66$\% \\
去除双层（MARL）& $\mathbf{74.98{\pm}45.99}$\% & $57.26{\pm}30.54$\% & $66.81{\pm}47.14$\% \\
\midrule
\textbf{完整 Bi-ARL} & $66.56{\pm}47.14$\% & $54.13{\pm}36.12$\% & $53.56{\pm}53.94$\% \\
\bottomrule
\end{tabular}
\end{table}

\subsection{对抗鲁棒性分析}
\label{sec:rob_cn}

本文在 NSL-KDD 和 UNSW-NB15 上同时采用 FGSM 与 PGD 对所有模型族进行鲁棒性评估。为控制正文篇幅，表~\ref{tab:rob_nsl_cn} 与表~\ref{tab:rob_unsw_cn} 仅保留 FGSM 的平均 FPR 汇总，完整的 FGSM/PGD 指标见结果文件 \texttt{robustness\_summary\_nsl\_kdd.csv} 与 \texttt{robustness\_summary\_unsw\_nb15.csv}。其中 NSL-KDD 上的 BiAT-MLP 鲁棒性结果仅为单种子冒烟测试，因此在表中单独标注。

\begin{table}[H]
\centering
\caption{NSL-KDD FGSM 鲁棒性（平均 FPR；RL 模型为 4 种子，BiAT-MLP 为单种子冒烟测试）}
\label{tab:rob_nsl_cn}
\begin{tabular}{lccc}
\toprule
\textbf{模型} & \textbf{干净} & \textbf{FGSM-$\epsilon$0.1} & \textbf{FGSM-$\epsilon$0.3} \\
\midrule
Bi-ARL        & 10.41\% & 47.03\% & 100.00\% \\
MARL          & 9.54\%  & 48.16\% & 99.81\% \\
Vanilla PPO   & 39.57\% & 54.06\% & 76.62\% \\
BiAT-MLP (1 seed) & 2.56\%  & 7.70\%  & 92.43\% \\
\bottomrule
\end{tabular}
\end{table}

\begin{table}[H]
\centering
\caption{UNSW-NB15 FGSM 鲁棒性（平均 FPR，4 个种子）}
\label{tab:rob_unsw_cn}
\begin{tabular}{lccc}
\toprule
\textbf{模型} & \textbf{干净} & \textbf{FGSM-$\epsilon$0.1} & \textbf{FGSM-$\epsilon$0.3} \\
\midrule
Bi-ARL               & 53.56\% & 77.85\% & 63.29\% \\
MARL                 & 66.81\% & 71.02\% & 80.75\% \\
Vanilla PPO          & 58.72\% & 56.15\% & 56.92\% \\
BiAT-MLP             & 28.31\% & 51.20\% & 95.06\% \\
Transformer-IDS（无AT）& \textbf{15.17\%} & 61.18\% & 70.86\% \\
\textbf{BiAT-FTTransformer} & 14.97\% & \textbf{47.68\%} & \textbf{54.56\%} \\
\bottomrule
\end{tabular}
\end{table}

核心发现：（1）NSL-KDD 上的 RL 模型在 FGSM-$\epsilon$0.3 下 FPR 接近 100\%，说明当前 RL 双层训练不足以应对强梯度攻击；（2）在 UNSW-NB15 上，Transformer-IDS（无对抗训练）与 BiAT-FTTransformer 的干净 FPR 相近（15.17\% 对 14.97\%），但在 FGSM-$\epsilon$0.3 下分别为 70.86\% 和 54.56\%，差距 16.3 个百分点。PGD 下同样呈现这一趋势：BiAT-FTTransformer 在 $\epsilon=0.1/0.3$ 时的 FPR 分别为 73.43\% / 96.56\%，而 BiAT-MLP 退化为 51.83\% / 100.00\%。这说明双层对抗训练的主要收益体现在对抗条件下的稳健性改善，而非干净数据上的绝对最优精度，与 TRADES 的鲁棒性-准确性权衡一致。

\subsection{超参数敏感性分析}

表~\ref{tab:kinner_cn} 报告 NSL-KDD 上 $K_{\text{inner}}$ 的扫描结果（各 2 个种子）。F1 在 $K_{\text{inner}}=5$ 时达到峰值（84.10\%），在 $K_{\text{inner}}=10$ 时略有下降（80.02\%，攻击者过度优化），验证了论文默认值 $K_{\text{inner}}=5$ 的合理性。

\begin{table}[H]
\centering
\caption{NSL-KDD：$K_{\text{inner}}$ 敏感性（均值 $\pm$ 标准差，2 个种子）}
\label{tab:kinner_cn}
\begin{tabular}{lccc}
\toprule
$K_{\text{inner}}$ & \textbf{F1} & \textbf{FPR} & \textbf{Recall} \\
\midrule
1  & $73.33{\pm}0.97$\% & $2.64{\pm}0.51$\%   & $59.05{\pm}1.46$\% \\
3  & $79.17{\pm}0.09$\% & $9.38{\pm}2.00$\%   & $70.17{\pm}1.12$\% \\
5  & $\mathbf{84.10{\pm}1.10}$\% & $21.13{\pm}11.13$\% & $84.23{\pm}8.01$\% \\
10 & $80.02{\pm}5.38$\% & $11.30{\pm}10.57$\% & $72.87{\pm}13.47$\% \\
\bottomrule
\end{tabular}
\end{table}

UNSW-NB15 上的 $\lambda$ 扫描（BiAT-FTTransformer 对抗损失权重）已完成定量统计。对 $\lambda \in \{0.3,0.5,0.7,0.9\}$ 的结果显示，随着 $\lambda$ 增大，F1 从 88.33\% 下降至 86.11\%，FPR 从 20.08\% 上升至 32.91\%。这表明在当前设置下，较小到中等的对抗损失权重更合适，而过大的对抗强调会破坏干净数据上的工作点。详细结果见 \texttt{sensitivity\_lambda\_unsw\_nb15.csv}。

\subsection{统计显著性}

表~\ref{tab:sig_cn} 报告 Wilcoxon 符号秩检验结果（$n=4$ 种子）。由于 $n=4$ 时检验功效不足，$p$ 值均未达到显著水平；但所有关键对比均表现出大效应量（Cohen's $d \geq 0.8$），提供了强有力的实际显著性证据。

\begin{table}[H]
\centering
\caption{关键对比统计显著性（$n=4$ 种子）}
\label{tab:sig_cn}
\begin{tabular}{llccc}
\toprule
\textbf{数据集} & \textbf{对比} & \textbf{指标} & \textbf{$\Delta$} & \textbf{Cohen's $d$} \\
\midrule
NSL-KDD   & Bi-ARL vs. Vanilla PPO & F1  & $+6.04$\% & 1.66（大效应） \\
UNSW-NB15 & BiAT-FTT vs. BiAT-MLP  & F1  & $+3.11$\% & 4.45（大效应） \\
UNSW-NB15 & BiAT-FTT vs. BiAT-MLP  & FPR & $-13.34$\% & $-4.84$（大效应） \\
UNSW-NB15 & BiAT-FTT vs. LightGBM  & FPR & $-2.17$\% & $-0.84$（大效应） \\
\bottomrule
\multicolumn{5}{l}{\scriptsize{$n{=}4$ 检验功效不足，以 Cohen's $d$ 作为补充证据。BiAT-FTT = BiAT-FTTransformer}}
\end{tabular}
\end{table}

\subsection{讨论}

综合全部实验，可以得出四个结论。（1）\textbf{Bi-ARL 建立了有原则的 RL 基线}：NSL-KDD 上的 Stackelberg 双层训练产生最佳 RL F1（80.17\%，Cohen's $d=1.66$），FPR 降低 29.16 个百分点。（2）\textbf{BiAT 框架的价值主要在对抗鲁棒性}：Transformer-IDS 与 BiAT-FTTransformer 在干净数据上性能相近（差距 $<1\%$ F1），但在 FGSM 攻击下 BiAT-FTTransformer 的 FPR 低 16.3 个百分点，证实双层对抗训练的核心价值在于对抗条件下的鲁棒性改善，与 TRADES 理论一致。（3）\textbf{树模型在总体 F1 上仍具优势}：LightGBM 领先 UNSW-NB15，并在 CIC-IDS2017 上接近饱和。（4）\textbf{内层循环稳定性是核心瓶颈}：UNSW-NB15 消融实验和 NSL-KDD 的 $\epsilon$=0.3 失效均指向内层循环不足以近似强攻击者这一工程问题。

%=============================================================================
\section{结论}
%=============================================================================

本文提出了双层对抗强化学习（Bi-ARL）框架及其监督式扩展——BiAT 框架，将攻防交互形式化为 Stackelberg 博弈，针对反复逼近最优响应的对手训练检测器，从而超越标准同步更新 MARL 提供了更原则化的鲁棒性训练目标。

在 NSL-KDD 上，Bi-ARL 以 80.17\% 的平均 F1 居强化学习基线之首，并将 FPR 相较 Vanilla PPO 降低 29.16 个百分点（Cohen's $d=1.66$，大效应量）。BiAT 框架将双层原则迁移至监督式检测器：BiAT-FTTransformer 在 UNSW-NB15 上取得 89.51\% 的 F1 与 14.97\% 的 FPR（所有神经网络检测器中最低），与 LightGBM 的 F1 差距仅 2.35 个点。新增的 Transformer-IDS 基线验证了 BiAT 框架的核心贡献：双层对抗训练以约 1\% 的干净 F1 换取了 FGSM-$\epsilon$0.3 下 16.3 个百分点的 FPR 降低（70.86\%$\to$54.56\%），与 TRADES 对抗训练理论一致。$K_{\text{inner}}$ 敏感性实验证实 $K_{\text{inner}}=5$ 是最优配置，效应量分析表明关键对比均具有大效应量。

本文将双层对抗训练确立为对抗鲁棒表格 IDS 的有效方法，明确了当前局限性并指出了具体的改进路径：（1）通过熵正则化或热身调度进一步稳定内层循环；（2）与近期基于自监督学习和图神经网络的 IDS 方法进行对比；（3）在 CSE-CIC-IDS2018、NF-UQ-NIDS 与 5G-NIDD 等更现代的数据集上进行评估，以超越经典基准的局限；（4）探索将双层对抗训练与 boosted-tree 集成方法相结合，以弥补当前在总体 F1 上的差距。

\bibliographystyle{plain}
\bibliography{../latex_source/references}

\end{document}
